\section{Discussion and Limitations}
\label{sec:discussion}
\textbf{Implicit visual consensus versus coordinate warping.} Bypassing externally supplied pose coordinates changes the dominant collaborative failure mode. Traditional intermediate feature fusion treats spatial alignment as an external prerequisite, using noisy transforms to warp dense feature grids. When misregistration occurs, spatial cues are injected into incorrect receptive fields, dispersing attention and creating duplicate ghost detections that cause negative collaborative gain. In contrast, \method{} delegates cross-agent spatial reasoning to multi-view token exchange within the shared visual backbone. Poses are inferred from visual evidence rather than supplied as a priori transforms, though residual error remains: the predicted-pose score is 14.53 percentage points below the GT-pose control at AP$_{.5}$.

\textbf{Limitations and open challenges.} Several open challenges remain before commercial deployment: (1)~\emph{Scale-aligned detection and pose error:} the reported AP removes global scale discrepancy at scoring time and does not by itself establish metric-scale detection; the 14.53-point GT-versus-predicted-pose gap remains. (2)~\emph{Communication streaming:} while JPEG compression reduces multi-camera payload by over 90\%, streaming raw images across dense fleets strains vehicular bandwidth; developing edge token pruning or latent feature streaming warrants investigation. (3)~\emph{Long-range scale sensitivity:} because box center errors grow with distance ($\Delta \mathbf{x} = \epsilon_s \mathbf{x}$), slight residual scale discrepancies degrade unaligned detection at long ranges; coupling multi-camera triangulation with temporal scale filtering remains an essential direction. (4)~\emph{Adverse conditions:} current cross-domain evidence does not cover heavy rain, dense fog, or packet dropouts, remaining critical frontiers for cooperative driving.
